% Stokes+isotropy => g=-M/A. Hop E_int is overlap, not 1/r.
\documentclass[11pt]{article}
\usepackage[margin=1.05in]{geometry}
\usepackage{amsmath,amssymb,amsthm}
\usepackage[colorlinks=true,linkcolor=blue,citecolor=blue,urlcolor=blue]{hyperref}

\newtheorem*{admission}{Admission}

\title{\textbf{The Flux Law Is $g=-M/A$.\\
The Hop Has No $1/r$.}}
\author{Robert W.\ McGwier\\
\small CTO, Cohere Technology Group\\
\small GitHub: \texttt{n4hy}}
\date{16 August 2026 --- Version 1}

\begin{document}
\maketitle

\begin{abstract}
\noindent
Entanglement supplies $M(B)=\sum_{i\in B}\delta e$.
If a local $\mathbf{g}$ exists,
\[
\oint_{\partial B}\mathbf{g}\cdot d\mathbf{A}=-M(B).
\]
Isotropy on a source-centered ball forces $g_R=-M/A$,
with $A$ the outgoing nearest-neighbour count.
That is Stokes, not a hop theorem.

Used: star slope $-1.9998$, sea $+0.992$.
Same class as $-M/R^2$ because $A/R^2\approx 19.4$.
Two-packet $E_{\mathrm{int}}$: $2.8\%$ at $d=2$
(overlap), $5\times 10^{-8}$ at $d=6$.
Successive $\lvert E_{\mathrm{int}}\rvert$ ratios
$5.6,16,48$: not $1/r$. Auditor: slopes confirmed;
$d=2$ overlap refutes a $2\%$ flat gate.

Gauss as used is the flux completion of the
measured mass. It is not in $H$. Not Einstein.
\end{abstract}

\begin{admission}
I asked $H$ for a force. At $d=2$ the packets
overlap and C\_flat failed. I did not move it.
The tail is not Coulomb.
\end{admission}

\end{document}
